\documentclass[12pt,a4paper]{article}
\usepackage[utf8]{inputenc}
\usepackage[T1]{fontenc}
\usepackage[indonesian]{babel}
\usepackage{geometry}
\usepackage{longtable}
\usepackage{array}
\usepackage{booktabs}
\usepackage{enumitem}
\usepackage[hidelinks,breaklinks]{hyperref}
\usepackage{xurl}
\geometry{margin=2.5cm}

\begin{document}

\begin{center}
\textbf{\Large RENCANA PEMBELAJARAN SEMESTER (RPS)}\\
\textbf{Berbasis Outcome-Based Education (OBE)}\\[0.5cm]
\textbf{PROGRAM STUDI TEKNIK INFORMATIKA}\\
\textbf{FAKULTAS TEKNIK}\\
\textbf{UNIVERSITAS LOREM IPSUM}\\[0.3cm]
\textbf{MATA KULIAH TEKNIK KOMPILASI}
\end{center}

\vspace{1cm}

% ============================================================
% 1. IDENTITAS MATA KULIAH
% ============================================================
\section*{1. Identitas Mata Kuliah}
\begin{tabular}{ll}
Nama Program Studi & : Teknik Informatika \\
Nama Mata Kuliah & : Teknik Kompilasi \\
Kode Mata Kuliah & : TIF-401 \\
Semester & : 6 (Enam) \\
SKS / Bobot Kredit & : 3 SKS (2 Teori, 1 Praktikum) \\
Dosen Pengampu & : Tim Dosen Teknik Informatika \\
Tanggal Penyusunan & : 5 Februari 2026 \\
Prasyarat & : Struktur Data, Algoritma, Pemrograman \\
\end{tabular}

\vspace{0.5cm}

% ============================================================
% 2. CAPAIAN PEMBELAJARAN LULUSAN (CPL)
% ============================================================
\section*{2. Capaian Pembelajaran Lulusan (CPL)}

CPL yang dibebankan pada mata kuliah Teknik Kompilasi ini mencakup dua capaian pembelajaran lulusan, yaitu:

\begin{itemize}[leftmargin=*]
  \item \textbf{CPL 02:} Memiliki pengetahuan yang memadai terkait dengan cara kerja sistem komputer dan mampu merancang dan mengembangkan berbagai algoritma metode untuk memecahkan masalah.

  \item \textbf{CPL 03:} Kemampuan merancang dan menganalisis algoritma untuk menyelesaikan permasalahan organisasi secara optimal, serta memilih dan menerapkannya pada bahasa pemrograman tertentu.
\end{itemize}

Kedua CPL tersebut dijabarkan menjadi Capaian Pembelajaran Mata Kuliah (CPMK) pada bagian berikut.

\vspace{0.5cm}

% ============================================================
% 3. CAPAIAN PEMBELAJARAN MATA KULIAH (CPMK)
% ============================================================
\section*{3. Capaian Pembelajaran Mata Kuliah (CPMK)}

Kemampuan atau kompetensi spesifik yang diharapkan mahasiswa kuasai setelah menyelesaikan mata kuliah:

\begin{itemize}[leftmargin=*]
  \item \textbf{CPMK 01:} Menjelaskan tahapan kompilasi (lexical, syntax, semantic analysis).

  \item \textbf{CPMK 02:} Memahami konsep Otomata Hingga (Finite Automata).

  \item \textbf{CPMK 03:} Merancang tata bahasa (grammar) bebas konteks untuk bahasa sederhana.

  \item \textbf{CPMK 04:} Mengimplementasikan parser sederhana untuk validasi kode.
\end{itemize}

\vspace{0.5cm}

% ============================================================
% 4. SUB-CPMK / INDIKATOR PENCAPAIAN
% ============================================================
\section*{4. Sub-CPMK / Indikator Pencapaian}

Penjabaran setiap CPMK menjadi indikator yang lebih terukur dan dapat diuji:

\noindent\textbf{CPMK 01: Tahapan Kompilasi (Leksikal, Sintaksis, Semantik)}
\begin{itemize}[leftmargin=*]
  \item \textbf{Sub-CPMK 1.1:} Menjelaskan perbedaan interpreter dan compiler serta jenis-jenis translator (assembler, interpreter, compiler)
  \item \textbf{Sub-CPMK 1.2:} Mengidentifikasi fase-fase kompilator (analisis leksikal, sintaksis, semantik) beserta peran \textit{symbol table management} dan \textit{error handling}
  \item \textbf{Sub-CPMK 1.3:} Menganalisis trade-off \textit{one-pass} vs \textit{multi-pass compiler} dan menjelaskan pengaruh \textit{run-time environment} (static, stack, heap allocation)
  \item \textbf{Sub-CPMK 1.4:} Menjelaskan hierarki Chomsky (Type 3--0) dan hubungannya dengan automata (FA, PDA, TM)
  \item \textbf{Sub-CPMK 1.5:} Menjelaskan pengertian analisis semantik, \textit{attribute grammar} (synthesized/inherited), serta \textit{syntax-directed definition} (SDD) dan \textit{translation schemes}
  \item \textbf{Sub-CPMK 1.6:} Mengimplementasikan \textit{symbol table} dengan \textit{nested scopes} dan melakukan \textit{type checking} untuk ekspresi dan statement
  \item \textbf{Sub-CPMK 1.7:} Menangani \textit{semantic error} dengan pesan yang informatif
\end{itemize}

\vspace{0.2cm}

\noindent\textbf{CPMK 02: Otomata Hingga (Finite Automata)}
\begin{itemize}[leftmargin=*]
  \item \textbf{Sub-CPMK 2.1:} Membuat regular expression untuk token spesifik bahasa
  \item \textbf{Sub-CPMK 2.2:} Mengimplementasikan NFA dan DFA untuk \textit{token recognition}
  \item \textbf{Sub-CPMK 2.3:} Menjelaskan ekivalensi NFA--DFA dan konsep minimalisasi DFA
  \item \textbf{Sub-CPMK 2.4:} Mengimplementasikan \textit{lexical analyzer} (\textit{hand-written} dan \textit{generator-based}) untuk bahasa sederhana
\end{itemize}

\vspace{0.2cm}

\noindent\textbf{CPMK 03: Tata Bahasa Bebas Konteks (Context-Free Grammar)}
\begin{itemize}[leftmargin=*]
  \item \textbf{Sub-CPMK 3.1:} Mendefinisikan CFG dengan notasi BNF dan EBNF, serta membuat derivasi dan \textit{parse tree}
  \item \textbf{Sub-CPMK 3.2:} Menganalisis ambiguitas grammar, \textit{left recursion}, dan \textit{left factoring}
  \item \textbf{Sub-CPMK 3.3:} Menghitung himpunan FIRST dan FOLLOW serta menjelaskan sifat grammar LL(1) dan LR(1)
  \item \textbf{Sub-CPMK 3.4:} Merancang grammar bebas konteks untuk bahasa sederhana (misalnya subset bahasa C)
\end{itemize}

\vspace{0.2cm}

\noindent\textbf{CPMK 04: Implementasi Parser Sederhana untuk Validasi Kode}
\begin{itemize}[leftmargin=*]
  \item \textbf{Sub-CPMK 4.1:} Membangun \textit{recursive descent parser} dan \textit{predictive parser} (LL(1))
  \item \textbf{Sub-CPMK 4.2:} Menerapkan teknik \textit{bottom-up parsing} (shift-reduce, operator precedence, SLR/LALR) serta menganalisis konflik shift-reduce dan reduce-reduce
  \item \textbf{Sub-CPMK 4.3:} Menggunakan \textit{parser generator} (Flex/Bison) dan \textit{semantic actions}
  \item \textbf{Sub-CPMK 4.4:} Membangun AST dan memvalidasi kode dengan pesan kesalahan sintaksis yang informatif
\end{itemize}

\vspace{0.5cm}

% ============================================================
% 5. VALIDASI SUB-CPMK TERHADAP REFERENSI
% ============================================================
\section*{5. Validasi Sub-CPMK terhadap Referensi}

Seluruh Sub-CPMK pada butir 4 divalidasi terhadap ketersediaan referensi terbuka yang dapat diunduh (PDF/HTML). Hasil validasi menunjukkan bahwa \textbf{seluruh 19 Sub-CPMK valid} dan dapat didukung oleh referensi pada butir 12; tidak ada Sub-CPMK yang dihapus.

\begin{longtable}{|>{\raggedright\arraybackslash}p{5.5cm}|>{\raggedright\arraybackslash}p{6cm}|}
\hline
\textbf{Sub-CPMK} & \textbf{Referensi Pendukung (butir 12)} \\
\hline
\endfirsthead
\hline
\textbf{Sub-CPMK} & \textbf{Referensi Pendukung (butir 12)} \\
\hline
\endhead
\hline
\endfoot

Sub-CPMK 1.1 & [1], [3], [4], [6] \\
\hline
Sub-CPMK 1.2 & [1], [3], [9], [11] \\
\hline
Sub-CPMK 1.3 & [1], [2], [4] \\
\hline
Sub-CPMK 1.4 & [5], [16] \\
\hline
Sub-CPMK 1.5 & [1], [6], [7] \\
\hline
Sub-CPMK 1.6 & [2], [3], [8] \\
\hline
Sub-CPMK 1.7 & [1], [3], [7] \\
\hline
Sub-CPMK 2.1 & [1], [2], [6] \\
\hline
Sub-CPMK 2.2 & [2], [6], [16] \\
\hline
Sub-CPMK 2.3 & [2], [5], [16] \\
\hline
Sub-CPMK 2.4 & [1], [3], [12], [15] \\
\hline
Sub-CPMK 3.1 & [1], [6], [16] \\
\hline
Sub-CPMK 3.2 & [1], [6], [10] \\
\hline
Sub-CPMK 3.3 & [6], [8], [10] \\
\hline
Sub-CPMK 3.4 & [1], [3], [8] \\
\hline
Sub-CPMK 4.1 & [1], [6], [10] \\
\hline
Sub-CPMK 4.2 & [2], [6], [8] \\
\hline
Sub-CPMK 4.3 & [12], [13], [14] \\
\hline
Sub-CPMK 4.4 & [1], [7], [15] \\
\hline
\end{longtable}

\vspace{0.5cm}

% ============================================================
% 6. MATERI PEMBELAJARAN (BAHAN KAJIAN)
% ============================================================
\section*{6. Materi Pembelajaran (Bahan Kajian)}

Empat belas materi pembelajaran disusun berdasarkan CPL, CPMK, dan Sub-CPMK, dengan rujukan pada referensi terbuka pada butir 12:

\begin{enumerate}[leftmargin=*]
  \item \textbf{Pengenalan Kompilasi dan Arsitektur Kompilator}\\
  Translator (assembler, interpreter, compiler); perbedaan compiler dan interpreter; fase \textit{front end} (analisis leksikal, sintaksis, semantik) dan \textit{back end}; \textit{symbol table management} dan \textit{error handling}; \textit{one-pass} vs \textit{multi-pass compiler}; \textit{run-time environment} (static, stack, heap allocation).

  \item \textbf{Teori Bahasa Formal dan Hierarki Chomsky}\\
  Alfabet, string, dan bahasa; grammar dan derivasi; hierarki Chomsky Type 3--0; hubungan kelas bahasa dengan automata (FA, PDA, TM).

  \item \textbf{Analisis Leksikal: Token, Lexeme, Pattern, dan Ekspresi Reguler}\\
  Token, lexeme, dan pattern; definisi, operasi, dan notasi ekspresi reguler; spesifikasi token dengan ekspresi reguler.

  \item \textbf{Finite Automata: NFA, DFA, Konversi, dan Minimalisasi}\\
  Definisi NFA dan DFA; ekivalensi NFA--DFA; konversi NFA ke DFA; minimalisasi DFA; penggunaan finite automata untuk \textit{token recognition}.

  \item \textbf{Implementasi Lexical Analyzer}\\
  Struktur data dan algoritma scanning; lexer \textit{hand-written}; lexer \textit{generator-based} (Flex); penanganan kesalahan leksikal.

  \item \textbf{Context-Free Grammar}\\
  Definisi CFG; notasi BNF dan EBNF; derivasi (leftmost/rightmost); parse tree dan pengantar abstract syntax tree.

  \item \textbf{Ambiguitas Grammar, Left Recursion, dan Left Factoring}\\
  Ambiguitas grammar dan cara mengatasinya; left recursion dan penghilangannya; left factoring; penggunaan precedence dan associativity.

  \item \textbf{Top-Down Parsing}\\
  Recursive descent parser; predictive parser; himpunan FIRST dan FOLLOW; sifat dan konstruksi parser LL(1); penanganan kesalahan sintaksis.

  \item \textbf{Bottom-Up Parsing}\\
  Shift-reduce parser; operator precedence parser; konsep handle dan item.

  \item \textbf{LR Parsing}\\
  Canonical LR, SLR, dan LALR; konstruksi tabel parsing; konflik shift-reduce dan reduce-reduce; kelas grammar LL(k) dan LR(k).

  \item \textbf{Parser Generator dan Semantic Actions}\\
  Flex dan Bison; ANTLR; semantic actions; integrasi lexer dan parser.

  \item \textbf{Analisis Semantik}\\
  Pengertian analisis semantik; attribute grammar (synthesized dan inherited); syntax-directed definition (SDD); translation schemes.

  \item \textbf{Symbol Table, Scope Management, dan Type Checking}\\
  Symbol table dan implementasinya; nested scopes; type system dan type expression; type checking ekspresi dan statement; type conversion dan type coercion.

  \item \textbf{Abstract Syntax Tree, Validasi Kode, dan Penanganan Error}\\
  Konstruksi dan traversal AST; validasi kode; jenis-jenis semantic error; penanganan dan pelaporan kesalahan yang informatif.
\end{enumerate}

\vspace{0.5cm}

% ============================================================
% 7. RENCANA PEMBELAJARAN MINGGUAN
% ============================================================
\section*{7. Rencana Pembelajaran Mingguan}

Dari 16 pertemuan dalam satu semester, 14 pertemuan digunakan untuk materi (butir 6) dan 2 pertemuan untuk UTS dan UAS:

\begin{longtable}{|c|>{\raggedright\arraybackslash}p{5.3cm}|>{\raggedright\arraybackslash}p{2.7cm}|>{\raggedright\arraybackslash}p{2.6cm}|}
\hline
\textbf{Pert.} & \textbf{Materi / Topik} & \textbf{Sub-CPMK} & \textbf{Referensi} \\
\hline
\endfirsthead
\hline
\textbf{Pert.} & \textbf{Materi / Topik} & \textbf{Sub-CPMK} & \textbf{Referensi} \\
\hline
\endhead
\hline
\endfoot

1 & Pengenalan Kompilasi dan Arsitektur Kompilator & 1.1, 1.2, 1.3 & [1], [3], [9], [11] \\
\hline
2 & Teori Bahasa Formal dan Hierarki Chomsky & 1.4 & [5], [16] \\
\hline
3 & Analisis Leksikal: Token, Lexeme, Pattern, dan Ekspresi Reguler & 2.1 & [1], [2], [6] \\
\hline
4 & Finite Automata: NFA, DFA, Konversi, dan Minimalisasi & 2.2, 2.3 & [2], [5], [16] \\
\hline
5 & Implementasi Lexical Analyzer (hand-written dan Flex) \textit{--- Praktikum 1} & 2.4 & [1], [3], [12], [15] \\
\hline
6 & Context-Free Grammar: BNF/EBNF, Derivasi, dan Parse Tree & 3.1 & [1], [6], [16] \\
\hline
7 & Ambiguitas Grammar, Left Recursion, dan Left Factoring & 3.2 & [1], [6], [10] \\
\hline
\textbf{8} & \textbf{Ujian Tengah Semester (UTS)} & \textbf{CPMK 01--03} & [1], [6], [8] \\
\hline
9 & Top-Down Parsing: Recursive Descent, LL(1), FIRST dan FOLLOW \textit{--- Praktikum 2} & 3.3, 4.1 & [1], [6], [10] \\
\hline
10 & Bottom-Up Parsing: Shift-Reduce dan Operator Precedence & 4.2 & [2], [6], [8] \\
\hline
11 & LR Parsing: Canonical LR, SLR, LALR, dan Konflik Parsing & 4.2 & [1], [6], [8] \\
\hline
12 & Parser Generator (Flex/Bison, ANTLR) dan Semantic Actions \textit{--- Praktikum 3} & 4.3 & [12], [13], [14] \\
\hline
13 & Analisis Semantik: Attribute Grammar, SDD, dan Translation Schemes & 1.5 & [1], [6], [7] \\
\hline
14 & Symbol Table, Scope Management, dan Type Checking \textit{--- Praktikum 4} & 1.6 & [2], [3], [8] \\
\hline
15 & Abstract Syntax Tree, Validasi Kode, dan Penanganan Error & 4.4, 1.7 & [1], [3], [4], [7] \\
\hline
\textbf{16} & \textbf{Ujian Akhir Semester (UAS)} & \textbf{CPMK 01--04} & [1], [6], [8] \\
\hline
\end{longtable}

\vspace{0.5cm}

% ============================================================
% 8. METODE PEMBELAJARAN
% ============================================================
\section*{8. Metode Pembelajaran}

Strategi atau pendekatan pembelajaran yang dipilih sesuai OBE yang menekankan aktivitas mahasiswa:

\begin{itemize}[leftmargin=*]
  \item \textbf{Ceramah Interaktif:} Penjelasan konsep kompilator dengan diskusi tanya jawab
  \item \textbf{Problem-Based Learning (PBL):} Mahasiswa menyelesaikan permasalahan kompilasi nyata
  \item \textbf{Project-Based Learning:} Pengembangan \textit{front-end} compiler secara bertahap sebagai proyek
  \item \textbf{Praktikum Terbimbing:} Implementasi komponen kompilator dengan bimbingan
  \item \textbf{Peer Review:} Mahasiswa melakukan code review terhadap implementasi compiler
  \item \textbf{Flipped Classroom:} Mahasiswa mempelajari materi teori sebelum kelas, praktik di kelas
  \item \textbf{Studi Kasus:} Analisis compiler nyata (GCC, Clang, TinyCC)
\end{itemize}

\vspace{0.5cm}

% ============================================================
% 9. PENGALAMAN BELAJAR MAHASISWA
% ============================================================
\section*{9. Pengalaman Belajar Mahasiswa}

Deskripsi tugas, aktivitas, atau pengalaman belajar yang mendukung pencapaian Sub-CPMK:

\begin{itemize}[leftmargin=*]
  \item Mengidentifikasi token, lexeme, dan pattern pada contoh bahasa sederhana
  \item Membuat derivasi dan parse tree untuk kalimat dari grammar yang diberikan
  \item Menghitung himpunan FIRST dan FOLLOW untuk grammar
  \item Menganalisis konflik shift-reduce dan reduce-reduce pada tabel LR
  \item Menerapkan attribute grammar (synthesized/inherited) pada pohon syntax
  \item Membaca dan merangkum konsep dari buku/referensi terbuka (Thain, Mogensen, Bergmann)
  \item Mengerjakan soal formal: desain grammar untuk bahasa sederhana, analisis ambiguitas, konstruksi NFA/DFA dari RE, penjelasan konflik parsing
  \item Presentasi singkat konsep teori (hierarki Chomsky, NFA vs DFA, LL vs LR, atau analisis semantik)
  \item Menganalisis trade-off teori (expressive power vs complexity) untuk pilihan desain kompilator
  \item Menganalisis arsitektur compiler sederhana (TinyCC, TCC)
  \item Mengimplementasikan lexer hand-written untuk subset bahasa C
  \item Membangun parser recursive descent untuk ekspresi aritmatika
  \item Menggunakan Flex dan Bison untuk parser generation
  \item Mengimplementasikan symbol table dengan nested scopes
  \item Melakukan type checking dan semantic analysis
  \item Berkolaborasi dalam tim untuk mengembangkan front-end compiler lengkap
  \item Mempresentasikan hasil proyek compiler
  \item Melakukan benchmarking dan analisis performa parser
\end{itemize}

\vspace{0.5cm}

% ============================================================
% 10. KRITERIA, INDIKATOR, DAN BOBOT PENILAIAN
% ============================================================
\section*{10. Kriteria, Indikator, dan Bobot Penilaian}

Teknik/alat asesmen dipetakan ke Sub-CPMK/CPMK dengan bobot yang jelas:

\begin{longtable}{|>{\raggedright\arraybackslash}p{2.5cm}|>{\raggedright\arraybackslash}p{4cm}|>{\raggedright\arraybackslash}p{5.5cm}|c|}
\hline
\textbf{Komponen} & \textbf{Teknik Asesmen} & \textbf{Indikator/CPMK} & \textbf{Bobot (\%)} \\
\hline
\endfirsthead
\hline
\textbf{Komponen} & \textbf{Teknik Asesmen} & \textbf{Indikator/CPMK} & \textbf{Bobot (\%)} \\
\hline
\endhead
\hline
\endfoot

Kuis & Multiple Choice \& Essay & Sub-CPMK 1.1, 1.2, 1.3 & 5 \\
\hline
Tugas Teori / Problem Set & Soal grammar, automata, parsing & Sub-CPMK 1.4, 2.1, 2.3, 3.2, 3.3 & 8 \\
\hline
Tugas Praktikum 1 (Pert. 5) & Lexer Implementation & Sub-CPMK 2.1, 2.2, 2.4 & 8 \\
\hline
Tugas Praktikum 2 (Pert. 9) & Parser Implementation & Sub-CPMK 3.1, 3.4, 4.1 & 8 \\
\hline
Tugas Praktikum 3 (Pert. 12) & Parser Generator \& AST & Sub-CPMK 4.2, 4.3, 4.4 & 7 \\
\hline
Tugas Praktikum 4 (Pert. 14) & Symbol Table \& Semantic Analysis & Sub-CPMK 1.5, 1.6, 1.7 & 8 \\
\hline
UTS (Pert. 8) & Written Exam \& Coding (konsep teori: bahasa formal, automata, sifat grammar, parsing) & CPMK 01, CPMK 02, CPMK 03 & 20 \\
\hline
Project Final & Mini Compiler Front-End & CPMK 02, CPMK 03, CPMK 04 & 20 \\
\hline
UAS (Pert. 16) & Comprehensive Exam (konsep teori \& arsitektur kompilator, analisis semantik) & CPMK 01, CPMK 02, CPMK 03, CPMK 04 & 16 \\
\hline
\textbf{Total} & & & \textbf{100} \\
\hline
\end{longtable}

\textbf{Kriteria Penilaian:}
\begin{itemize}[leftmargin=*]
  \item A (85-100): Menguasai semua CPMK dengan sangat baik, mampu menerapkan dalam kasus kompleks
  \item B (70-84): Menguasai sebagian besar CPMK dengan baik
  \item C (60-69): Menguasai CPMK dasar dengan cukup
  \item D (50-59): Menguasai sebagian kecil CPMK
  \item E (<50): Belum menguasai CPMK yang ditetapkan
\end{itemize}

\textbf{Rubrik Penilaian Tugas Praktikum:}
\begin{itemize}[leftmargin=*]
  \item \textbf{Correctness (40\%):} Program berjalan dengan benar, mengatasi edge cases
  \item \textbf{Code Quality (30\%):} Struktur kode, readability, dokumentasi
  \item \textbf{Completeness (20\%):} Semua requirement terpenuhi
  \item \textbf{Testing (10\%):} Test cases yang komprehensif
\end{itemize}

\vspace{0.5cm}

% ============================================================
% 11. EVALUASI DAN REFLEKSI PEMBELAJARAN
% ============================================================
\section*{11. Evaluasi dan Refleksi Pembelajaran}

Penilaian sumatif/formatif untuk memantau ketercapaian outcome secara menyeluruh:

\begin{itemize}[leftmargin=*]
  \item \textbf{Evaluasi Formatif:} Kuis mingguan, progress review proyek compiler, peer review untuk memberikan feedback berkelanjutan
  \item \textbf{Evaluasi Sumatif:} UTS dan UAS untuk mengukur pencapaian CPMK secara komprehensif
  \item \textbf{Refleksi Mahasiswa:} Jurnal belajar mingguan untuk refleksi diri terhadap pemahaman materi kompilator
  \item \textbf{Evaluasi Dosen:} Survey kepuasan mahasiswa di tengah dan akhir semester
  \item \textbf{Continuous Improvement:} Analisis hasil penilaian untuk perbaikan RPS di semester berikutnya
\end{itemize}

\vspace{0.5cm}

% ============================================================
% 12. DAFTAR PUSTAKA
% ============================================================
\section*{12. Daftar Pustaka}

Seluruh referensi berikut merupakan sumber terbuka (\textit{open access}) yang dapat diunduh secara gratis dalam format PDF atau HTML, dan telah diverifikasi dapat diakses pada 9 Oktober 2026.

\noindent\textbf{A. Buku Teks Terbuka (PDF)}

\begin{enumerate}[leftmargin=*]
  \item Thain, D. (2024). \textit{Introduction to Compilers and Language Design} (2nd ed.). University of Notre Dame. PDF: \url{https://dthain.github.io/books/compiler/compilerbook.pdf}

  \item Mogensen, T. {\AE}. (2010). \textit{Basics of Compiler Design} (Anniversary ed.). University of Copenhagen. PDF: \url{https://hjemmesider.diku.dk/~torbenm/Basics/basics_lulu2.pdf} (solusi latihan: \url{https://hjemmesider.diku.dk/~torbenm/Basics/solutions.pdf})

  \item Bergmann, S. D. (2024). \textit{Compiler Design: Theory, Tools, and Examples}. Rowan University. PDF: \url{https://protect.bju.edu/cps/docs/cps450/BergmannCompilerDesign.pdf}

  \item Bornat, R. (2008). \textit{Understanding and Writing Compilers: A Do-It-Yourself Guide} (Internet ed.). Middlesex University. PDF: \url{https://www.eis.mdx.ac.uk/staffpages/r_bornat/books/compiling.pdf}

  \item Critchlow, C., \& Eck, D. (2024). \textit{Foundations of Computation} (v2.3.2). Open Textbook Library. PDF: \url{https://math.hws.edu/FoundationsOfComputation/FoundationsOfComputation_2.3.2_8.5x11.pdf}
\end{enumerate}

\noindent\textbf{B. Catatan Kuliah dan Slide Perkuliahan (PDF)}

\begin{enumerate}[leftmargin=*,start=6]
  \item Srikant, Y. N. \textit{Principles of Compiler Design --- NPTEL Lecture Notes}. Indian Institute of Science Bangalore. PDF: \url{https://nitsri.asia/Department/Computer%20Science%20&%20Engineering/CD-LEC-NOTES.pdf}

  \item MIT OpenCourseWare. (2010). \textit{6.035 Computer Language Engineering --- Lecture Notes}. Massachusetts Institute of Technology. \url{https://ocw.mit.edu/courses/6-035-computer-language-engineering-spring-2010/pages/lecture-notes/}

  \item Stanford University. \textit{CS143: Compilers --- Lecture Slides dan Handouts}. \url{https://web.stanford.edu/class/cs143/}; catatan parsing: Dill, D. L. \textit{CS143 Notes: Parsing}. \url{https://web.stanford.edu/class/archive/cs/cs143/cs143.1156/handouts/parsing.pdf}

  \item Universit{\"a}t des Saarlandes. \textit{The Structure of Compilers} (Compiler Construction slides). PDF: \url{https://compilers.cs.uni-saarland.de/teaching/cc/2013/slides/structure.pdf}

  \item McGill University. \textit{CS 520: Compiler Design --- Top-Down Parsing} handout. PDF: \url{https://cs.mcgill.ca/~cs520/2026/slides/parsing-handout.pdf}

  \item University of Edinburgh. \textit{Compiling Techniques --- The Structure of Compilers} (Open Course slides). PDF: \url{https://opencourse.inf.ed.ac.uk/sites/default/files/2025-01/ct-lecture2-35000feet.pdf}
\end{enumerate}

\noindent\textbf{C. Dokumentasi Alat dan Referensi Teori Bahasa}

\begin{enumerate}[leftmargin=*,start=12]
  \item Estes, W., Paxson, V., \& Millaway, J. \textit{Lexical Analysis With Flex} (Flex 2.6.2 Manual). \url{https://westes.github.io/flex/manual/}

  \item GNU Project. \textit{GNU Bison Manual --- The Yacc-Compatible Parser Generator}. \url{https://www.gnu.org/software/bison/manual/}

  \item ANTLR Project. \textit{ANTLR 4 Documentation}. \url{https://www.antlr.org/}

  \item Wikibooks. \textit{Compiler Construction} (open book). \url{https://en.wikibooks.org/wiki/Compiler_Construction}

  \item Wikibooks. \textit{Theory of Formal Languages, Automata, and Computation} (open book). \url{https://en.wikibooks.org/wiki/Theory_of_Formal_Languages,_Automata,_and_Computation}
\end{enumerate}

\vspace{0.3cm}

\noindent\textbf{Catatan:} Nomor rujukan [1]--[16] digunakan pada butir 5, 6, dan 7. Seluruh berkas dapat diunduh tanpa biaya; sebagian buku teks membatasi penyebaran ulang untuk kepentingan komersial, namun bebas digunakan untuk keperluan pembelajaran dan akademik.

\vspace{1cm}

\begin{flushright}
\begin{tabular}{c}
Disusun oleh,\\[2cm]
\textbf{Tim Dosen Teknik Informatika}\\
Program Studi S1 Teknik Informatika\\
Fakultas Teknik\\[0.5cm]
5 Februari 2026
\end{tabular}
\end{flushright}

\end{document}
